\documentclass[11pt]{article}
\usepackage[a4paper,margin=25mm]{geometry}
\usepackage{amsmath,amssymb,amsthm}
\usepackage[hidelinks]{hyperref}
\emergencystretch=2em
\newtheorem{theorem}{Theorem}
\newcommand{\Aut}{\operatorname{Aut}}
\title{The type-$A_2$ exchange graph contradicts\\the proposed Weyl semidirect product\\
\large Disproof of Conjecture 00000000565}
\author{gaochengzhecpu}\date{}
\begin{document}\maketitle
\begin{abstract}
The coefficient-free cluster exchange graph of type $A_2$ is a pentagon.
Its automorphism group has order $10$, while the Weyl group of $A_2$
has order $6$. Consequently its automorphism group cannot be
$W(A_2)\rtimes H$ for any group $H$ and any action. We derive the
exchange graph directly from symbolic cluster mutation, and the
accompanying Lean project verifies the complete counterexample.
\end{abstract}

\section*{Interpretation and the example}
We use the standard unlabelled cluster exchange graph: vertices are
clusters and edges are single mutations. This is the customary
finite/affine-type meaning of ``exchange graph'' in the conjecture.
The source does not give a different formal definition; in particular
we do not identify this graph with a Weyl chamber graph. The type-$A_2$
pentagon is classical; see Fomin--Zelevinsky, Example~7.6 and Figure~3.

Work in $K=\mathbb Q(x,y)$, where $x,y$ are independent indeterminates,
with initial exchange matrix
\[
 B=\begin{pmatrix}0&1\\-1&0\end{pmatrix}.
\]
Its Cartan companion is the Cartan matrix of $A_2$. In rank two each
mutation replaces $B$ by $-B$, so all exchange polynomials are $1+z$.
Thus mutating $a$ in the cluster $\{a,b\}$ gives
$\{(1+b)/a,b\}$. Define, with subscripts modulo $5$,
\[
 z_0=x,\quad z_1=y,\quad z_2=\frac{1+y}{x},\quad
 z_3=\frac{1+x+y}{xy},\quad z_4=\frac{1+x}{y}.
\]
All five are nonzero and pairwise distinct. For example, their
denominators are nonzero polynomials, and evaluation of these
fractions at $(x,y)=(2,5)$ gives
$2,5,3,4/5,3/5$, respectively. Equality of two rational functions
would force equality after this well-defined specialization.
Specialization is used only to certify distinctness, not to replace
the symbolic mutation problem by a numerical one.

\newpage
\section*{All clusters and their automorphisms}
Direct multiplication proves the five identities
\[
 z_i z_{i+2}=1+z_{i+1}\qquad(i\in\mathbb Z/5\mathbb Z).
\]
Set $C_i=\{z_i,z_{i+1}\}$. Their elements are distinct, and the five
sets $C_i$ are distinct. In $C_i$, mutation of $z_i$ produces
$C_{i+1}$, while mutation of $z_{i+1}$ produces $C_{i-1}$.
These are the only choices. Induction on a mutation sequence shows
that every cluster reachable from $C_0$ is one of the $C_i$.
Conversely, repeated forward mutation reaches all five. Hence the
exchange graph is exactly the cycle
\[
 C_0-C_1-C_2-C_3-C_4-C_0.
\]
An automorphism of a pentagon is determined by the image of one
vertex (five choices) and the image of a specified neighbor (two
choices); the remaining images are forced around the cycle.
All ten choices occur, via $i\mapsto b+i$ and $i\mapsto b-i$ modulo
$5$. Therefore
\[
 |\Aut(\mathcal E_{A_2})|=10.
\]

\section*{The Weyl group obstruction}
Realize $A_2$ using the roots $e_i-e_j$ in the plane
$v_0+v_1+v_2=0$. The reflection in $\alpha=e_i-e_j$ is
\[
 v\longmapsto v-\frac{2(v\cdot\alpha)}{\alpha\cdot\alpha}\alpha
   =v-(v_i-v_j)(e_i-e_j).
\]
It swaps coordinates $i,j$. The simple reflections
$s_0=(0\ 1)$ and $s_1=(1\ 2)$ generate $S_3$; its six elements are
\[
 1,\ s_0,\ s_1,\ s_0s_1,\ s_1s_0,\ s_0s_1s_0.
\]
The action on the root plane is faithful: a permutation fixing all
roots $e_i-e_j$ fixes their coordinate positions. Thus
$W(A_2)\cong S_3$ and $|W(A_2)|=6$.

\begin{theorem}
For every group $H$ and every action of $H$ on $W(A_2)$,
\[
 \Aut(\mathcal E_{A_2})\not\cong W(A_2)\rtimes H.
\]
\end{theorem}
\begin{proof}
Every such semidirect product contains an injective copy of
$W(A_2)$. An isomorphism would embed a group of order $6$ into a
group of order $10$, contradicting Lagrange's theorem.
\end{proof}
This also rules out the proposed graph-automorphism factor, whatever
action is chosen. A false finite-type statement cannot acquire a
valid affine extension merely by extending its scope. No claim about
a separately corrected affine conjecture is needed for this disproof.

\newpage
\section*{What the Lean project proves}
The ambient field is the actual fraction field of the two-variable
polynomial ring over $\mathbb Q$, using Mathlib's
\texttt{MvPolynomial} and \texttt{FractionRing} types.
\texttt{clusterVar\_injective} proves symbolic distinctness by
cross-multiplying in this field, using the injective polynomial
embedding, and only then applying polynomial evaluation.
\texttt{exchange\_relation} verifies the five rational identities.

\texttt{Mutates} is the two-variable $A_2$ rule stated above.
\texttt{mutation\_classification} proves that each $C_i$ has exactly
the two indicated possible successors. \texttt{Reachable} is an
inductive definition allowing arbitrary finite mutation sequences;
\texttt{reachable\_iff} proves that its states are exactly the five
listed clusters. An explicit equivalence transfers permutations of
these actual states to permutations of the five cycle indices.

\texttt{exchangeAut} is the subgroup of all permutations of the
reachable states preserving mutation in both directions.
\texttt{exchange\_aut\_card} proves its order is $10$. The finite
enumeration involved is checked by Lean's kernel.
\texttt{root\_length}, \texttt{root\_pairing}, and
\texttt{reflection\_is\_swap} verify the root reflection formula.
\texttt{WeylA2} is the subgroup generated by the two coordinate swaps;
\texttt{weyl\_card} proves its order is $6$ by exhaustive enumeration
of the six words above. Finally
\texttt{no\_semidirect\_description} excludes every second factor
$H$, even without a finiteness assumption on $H$.

The formalization specializes the standard seed-mutation definition
to its explicit coefficient-free $A_2$ rule. It does not claim to
formalize a general cluster-algebra classification or a general
root-system library. The connection to the source's standard
exchange-graph terminology is spelled out in the preceding proof.

\section*{Reproduction and review}
Lean 4.19.0 and Mathlib revision\\
\texttt{c44e0c8ee63ca166450922a373c7409c5d26b00b} are pinned.
In \texttt{lean/}, run
\begin{verbatim}
lake exe cache get
lake build
lake env lean -DwarningAsError=true Main.lean
\end{verbatim}
The project is freshly rebuilt against verified official dependencies.
It has no proof placeholders, custom axioms, external computational
oracle, or auxiliary numerical program. The dependency audit permits
only Lean's standard axioms. Run \texttt{tectonic main.tex} to regenerate
the PDF. A separate adversarial self-review was performed by the author;
no independent reviewer or subagent was used.

\section*{References}
\begin{enumerate}
\item The Last Math Competition, Conjecture 00000000565. Exact bilingual
source and upstream provenance are included.
\item S. Fomin and A. Zelevinsky, \emph{Cluster Algebras I: Foundations},
J. Amer. Math. Soc. 15 (2002), 497--529, Example~7.6 and Figure~3.
\href{https://www.esi.ac.at/preprints/esi1023.pdf}{Authors' preprint}.
\end{enumerate}
\end{document}
